\section{引言与相关工作}
网络流量矩阵描述一组源宿对在各时刻的流量。
当测量资源受限、部分监测不可用或观测在流间分配不均时，
补全器需要从有限可见值恢复缺失项。
同一条流的邻近观测通常提供直接的幅度和变化线索，
相关流的可见数据则可能补充本流在较长缺口中的信息。
如何组合这两类信息，并判断新增结构是否值得其计算成本，是本文的研究问题。

已有方法已广泛利用时空关系。
SPIN通过面向稀疏观测的时空信息传播，为目标位置构造可解码表示\cite{marisca2022spin}。
它本身能够读取本流与邻流的时间信息，因此不能把引入另一条记忆路径解释为
首次为该编码器增加时间依赖。
网络流量矩阵研究中的SRMF及其局部补正也表明，
较大范围估计与邻近真实观测相结合具有直接先例\cite{zhang2009srmf}。
本文采用局部观测直读的具体共享神经读出实现，不将全局与局部结合这一思想本身作为新颖性主张。

另一类相关结构是以矩阵状态存储键值关联的记忆。
Delta规则根据当前键的预测残差修正状态，DeltaNet研究了这一类更新及其高效训练\cite{yang2024deltanet}。
本文沿用项目中的同步桶更新实现：同一时刻的可见邻流从同一个旧状态计算修正，
以避免同桶流输入顺序直接决定更新结果。
这里不复用DeltaNet的语言建模实验结论，也没有实现其整套序列并行训练算法；
本文关注的是矩阵记忆作为流量补全分支时的实际价值。

我们构建共享上下文与解码接口的两个变体。
完整变体\full 将一层SPIN生成的上下文投影为记忆键、值和查询，
同时通过Direct路径读取本流前后各两个最近的真实观测。
去记忆变体\lite 保留上下文、Direct、查询与解码器，删除K/V和扫描。
两者全部参数从头共同训练，共有参数在同一初始化种子下取相同初值，
以较小的实验范围回答额外记忆是否提供可重复增量。

本文的具体工作包括：一是给出可见信息约束下的联合读出实现及可对应的去记忆变体；
二是在两套WAN快照、三种固定预算缺失条件和两个初始化下完整报告误差取舍；
三是在同一设备上测量包含编码、读出及输入输出处理的推理成本。
结果使后续候选收窄到较简单的\lite，
而\full 在平方误差指标上的收益及代价也被保留。
本文是一项开发性结构研究，目标是形成可复核的下一步判断，
并不宣称所有组成分支已被证明必要或当前方案普遍领先。
